%% =====================================================================
%%  Corrigés — Unit 11 « Can I have a coffee, please? »
%% =====================================================================
\corrigesUnite{11}{Can I have a coffee, please?}

\begin{corrige}{1}
\begin{reponses}[2]
\item \rep{It's fourteen dollars.}
\item \rep{It's sixteen dollars.}
\item \rep{It's nineteen dollars.}
\item \rep{They're six dollars.}
\item \rep{Yes, it is.}
\end{reponses}
\pourquoi{deux cafés à 3 dollars : 6 dollars. Pour deux choses, on répond avec \en{They're…}.}
\end{corrige}

\begin{corrige}{2}
\begin{reponses}[3]
\item \rep{I'd like}
\item \rep{I like}
\item \rep{I'd like}
\item \rep{I like}
\item \rep{I'd like}
\item \rep{I like}
\end{reponses}
\pourquoi{une demande (« je voudrais ») → \en{I'd like} ; un goût en général (« j'aime ») → \en{I like}.}
\end{corrige}

\begin{corrige}{3}
\begin{reponses}[5]
\item \rep{e}
\item \rep{a}
\item \rep{b}
\item \rep{c}
\item \rep{d}
\end{reponses}
\end{corrige}

\begin{corrige}{4}
\begin{reponses}[3]
\item \rep{can}
\item \rep{have}
\item \rep{here}
\item \rep{go}
\item \rep{I'd}
\item \rep{Would}
\end{reponses}
\end{corrige}

\begin{corrige}{5}
\begin{reponses}[1]
\item \rep{Can I have a coffee, please?}
\item \rep{I'd like the fish, please.}
\item \rep{Can we have the check, please?}
\item \rep{Could I have a glass of water, please?}
\end{reponses}
\pourquoi{les formules sont interchangeables : \en{Could I have the fish, please?} ou \en{I'd like a
coffee, please} sont aussi corrects. L'essentiel : une formule polie et \en{please}.}
\end{corrige}

\begin{corrige}{6}
\begin{reponses}[5][french]
\item \rep{b}
\item \rep{e}
\item \rep{a}
\item \rep{d}
\item \rep{c}
\end{reponses}
\end{corrige}

\begin{corrige}{7}
\begin{reponses}[2]
\item \rep{Can I have}
\item \rep{Could I have}
\item \rep{Would you like}
\item \rep{I'd like}
\item \rep{Can we have}
\item \rep{I'll have}
\end{reponses}
\pourquoi{texte de la piste page~\pageref{trans-11.10}. Remarquez comme \en{can}, \en{could} et
\en{would} se réduisent à \phon{kən}, \phon{kəd}, \phon{wəd}.}
\end{corrige}

\begin{corrige}{8}
\begin{reponses}[3]
\item \rep{tea}
\item \rep{small}
\item \rep{for here}
\item \rep{coffee}
\item \rep{large}
\item \rep{to go}
\item \rep{latte}
\item \rep{medium}
\item \rep{to go}
\end{reponses}
\pourquoi{Maya appelle Sam \en{Sam} et il lui répond \en{little sister} : c'est sa petite sœur.}
\end{corrige}

\begin{corrige}{9}
\begin{reponses}[4]
\item \rep{ready}
\item \rep{I'd}
\item \rep{Would}
\item \rep{please}
\item \rep{glass}
\item \rep{dessert}
\item \rep{thank}
\end{reponses}
\end{corrige}

\begin{corrige}{10}
\begin{reponses}[1]
\item \rep{I'd like a tea, please.}
\item \rep{For here or to go? — To go.}
\item \rep{Would you like something to drink? — Yes, please.}
\item \rep{Could I have the check, please?}
\item \rep{I'll have the burger with fries.}
\item \rep{The tip isn't included.}
\item \rep{I like coffee, but I'd like a tea today.}
\item \rep{It's delicious!}
\end{reponses}
\piege{phrase 7 : les deux verbes ont des sens différents. \en{I like} = j'aime ; \en{I'd like} = je
voudrais.}
\end{corrige}

\begin{corrige}{11}
Réponse modèle :
\begin{anglais}\otTableau
\begin{tabular}{@{}l@{\hspace{6mm}}l@{}}
\otNomDit{serveur} & Good evening! Are you ready to order?\\
\otNomDit{vous}    & Yes. I'd like the soup of the day, please.\\
\otNomDit{jack}    & And I'll have the burger and fries.\\
\otNomDit{serveur} & Would you like something to drink?\\
\otNomDit{vous}    & Two glasses of water, please.\\
\otNomDit{vous}    & Excuse me, can we have the check, please?\\
\otNomDit{serveur} & Sure. That's twenty-two dollars.\\
\end{tabular}
\end{anglais}
\retiens{6 + 16 = 22 dollars. Avec la taxe et un pourboire de 20 \%, comptez environ 28 dollars.}
\end{corrige}

\begin{corrige}{12}
Réponses modèles :
\begin{reponses}[1]
\item \rep{Can I have a coffee, please?}
\item \rep{Medium, please.}
\item \rep{To go, please.}
\item \rep{No, thank you.} \textit{ou} \rep{Yes, please!}
\item \rep{Here you are. Thank you!}
\end{reponses}
\end{corrige}

\begin{corrigetest}
\begin{reponses}[5][french]
\item \rep{c}
\item \rep{a}
\item \rep{b}
\item \rep{a}
\item \rep{c}
\item \rep{b}
\item \rep{a}
\item \rep{c}
\item \rep{b}
\item \rep{a}
\end{reponses}
\pourquoi{question 3 : \en{Yes, please} est la réponse polie et naturelle ; un simple \en{Yes} paraît
sec.}
\end{corrigetest}
